% Part: computability
% Chapter: recursive-functions
% Section: halting-problem

\documentclass[../../../include/open-logic-section]{subfiles}

\begin{document}

\olfileid{cmp}{rec}{hlt}
\olsection{مشكلة التوقف}

\emph{مشكلة التوقف} هي، على العموم، أن نبت هل ينتهي حساب دالة عند
الدخل~$n$ إلى نتيجة أم لا، متى أعطينا مواصفتها~$e$،
كبرنامجها مثلًا، وأعطينا العدد~$n$. ومن نتائج آلان تورنغ
الشهيرة أن هذا السؤال العام لا تجيب عنه دالة قابلة للحساب؛
فالدالة
\[
h(e, n) =
\begin{cases}
1 & \text{إذا توقف الحساب $e$ عند الدخل $n$}\\
0 & \text{في غير ذلك،}
\end{cases}
\]
ليست قابلة للحساب.

وأما في الدوال العودية الجزئية، فالفهرس~$e$ الوارد في
مبرهنة كليني للصورة السوية يقوم مقام مواصفة البرنامج. فمتى كانت $f$ دالة
عودية جزئية، عُدّ كل~$e$ يحقق معادلة تلك المبرهنة فهرسًا
لـ~$f$. ولكل عدد~$e$ تقضي مبرهنة الصورة السوية بأن
\[
\cfind{e}(x) \simeq U(\mu s \; T(e, x, s))
\]
دالة عودية جزئية، وأنه لكل دالة عودية جزئية
$f\colon\Nat\to\Nat$ يوجد $e\in\Nat$ بحيث
$\cfind{e}(x)\simeq f(x)$ لكل~$x\in\Nat$. بل لكل
$f$ من هذا القبيل ما لا يتناهى من الفهارس~$e$. وتعريف
\emph{دالة التوقف}~$h$ هو
\[
h(e, x) =
\begin{cases}
1 & \text{إذا كانت $\cfind{e}(x) \fdefined$}\\
0 & \text{في غير ذلك.}
\end{cases}
\]
وتكون $h(e,x)=0$ إذا وفقط إذا كانت
$\cfind{e}(x)\fundefined$.

% تصحيح مصدر (C226-01): كل عدد طبيعي فهرس في التعداد المعتمد، فحُذف الفرع المستحيل.
\textbf{تنبيه تصحيحي على الأصل.} يحصي الترميز المعتمد هنا دالة
عودية جزئية~$\cfind{e}$ لكل عدد طبيعي~$e$؛ فحُذف الفرع الذي يفترض
وجود عدد طبيعي لا يكون فهرسًا أصلًا.

\begin{thm}
\ollabel{thm:halting-problem}
دالة التوقف~$h$ ليست عودية جزئية.
\end{thm}

\begin{proof}
لنفرض عودية $h$ الجزئية؛ فيجوز حينئذ أن نعرّف
\[
d(y) =
\begin{cases}
1 & \text{إذا كان $h(y, y) = 0$}\\
\umin{x}{x \neq x} & \text{في غير ذلك.}
\end{cases}
\]
ولا عدد $x$ يحقق $x\neq x$؛ ولذلك لا تتعين
$\umin{x}{x\neq x}$، ومن ثم تكون $d(y)\fundefined$ إذا وفقط إذا كان
$h(y,y)\neq0$. فنحصل من التعريف على الأمرين الآتيين:
\begin{enumerate}
\item $d(y)\fdefined$ إذا وفقط إذا كان $h(y,y)=0$، أي إذا وفقط إذا
  كانت $\cfind{y}(y)\fundefined$.
\item $d(y)\fundefined$ إذا وفقط إذا كان $h(y,y)=1$، أي إذا وفقط إذا
  كانت $\cfind{y}(y)\fdefined$.
\end{enumerate}
ومن عودية $h$ الجزئية تلزم عودية $d$ الجزئية؛ فتجعل لها
مبرهنة كليني للصورة السوية فهرسًا~$e_d$، أي
$\cfind{e_d}\simeq d$. ولكن التعادلين السابقين يعطيان
\[
d(e_d)\fdefined \quad\text{إذا وفقط إذا}\quad
\cfind{e_d}(e_d)\fundefined
\quad\text{إذا وفقط إذا}\quad d(e_d)\fundefined\text{،}
\]
وهو محال. فبطل فرض عودية~$h$ الجزئية.
\end{proof}

\end{document}
